\section{Results}
\label{sec:results}

% Results report numbers. No mechanism, no conservation implication, no
% "suggesting that". Every one of those belongs in the Discussion.

\subsection{Connectivity change, \studyperiod}
% Sentence starters: "Across the study period, modelled connectivity ..." / "The largest changes occurred ..." / "At the core-pair level, ..."
\TODO{H1. Magnitude and geography of change.
PRIMARY MEASURE: effective resistance. Least-cost path cost is secondary and gives
bigger numbers for the same links; say which is which.
FACTS: 253 pairs, median -1.224 percent, median absolute 2.822, 90th percentile 6.089,
92 up and 161 down. The 45 selected links: median -1.557. The 32 priority links:
median -2.618, 12 up and 20 down. Biggest single move +13.1 percent on 33-38.
Outside-core current: median relative change 4.3 percent, top-decile agreement between
years 0.91 to 0.93.
TRAP: "report per country" cannot be done. Cores have no country attribution and the
only per-country table is covariate change, not connectivity. Either build the
attribution or drop the promise.
Figure~\ref{fig:change}.}

\begin{figure}[htbp]\centering
  \includegraphics[width=\textwidth]{Figure_0_study_area_and_core_groups.pdf}
  \caption{Study extent, fixed historical range cores, and core groups used to define
  the modeled connectivity network. Core groups are an analysis construct and do not
  imply separate populations.}
  \label{fig:studyarea}
\end{figure}

\subsection{Robust priority links}
% Sentence starters: "I identified [number] robust priority links ..." / "These links were concentrated ..." / "Their agreement across scenarios was ..."
\TODO{H3. How many links met the criterion and how concentrated they are.
FACTS: 32 of 45 primary, 13 uncertainty. Weight screen passes 34, fence screen removes
2 more. Vegetation-direction removes none that were still in.
NOTE: H3 predicted a LIMITED subset. 32 of 45 is most of the network, so it is only
partly supported. Worth saying, since it is a more interesting result than a clean
confirmation.
Table~\ref{tab:top20}.}

\begin{figure}[htbp]\centering
  \includegraphics[width=\textwidth]{Figure_5_modelled_links_and_transportation.pdf}
  \caption{Modeled core-to-core links with transportation and available fence context.
  Lines represent network links used to define least-cost analyses; they are not
  observed animal paths.}
  \label{fig:networkcontext}
\end{figure}

\subsection{Protected-area coverage of priority connectivity}
% Sentence starters: "Of the priority connectivity area, [value] was located ..." / "Compared with the predeclared null, this represented ..." / "The permutation test returned ..."
\TODO{H2. Protection of the connecting landscape. Rewrite, do not fill in.
WHY: the predeclared comparison included the fixed cores, and the whole top decile of
current sits inside them, so it measured a footprint that cannot change.
FACTS, outside cores, background 31.26 percent: 80th percentile 23.27, 90th 25.54,
95th 33.45, 99th 44.63. Crossover near the 93rd. Gap widened from -7.03 to -7.99
points at the 80th between 2012 and 2024. Direction holds at 0, 1, 5 and 10 km core
buffers.
TRAP: no p-value. It swings 0.006 to 0.096 on block size alone. Also drop the word
"observed"; this is modelled current.
Table~\ref{tab:pa}.}

\subsection{Sensitivity and plausibility checks}
% Sentence starters: "Results were consistent across ..." / "The links most sensitive to assumptions were ..." / "The plausibility assessment showed ..."
\TODO{Model sensitivity. Still substantial, but not the subsection that was planned.
DOES NOT EXIST: the cell-level agreement grid, and any agreement raster or map.
WHAT TO REPORT INSTEAD, and this is the paper's strongest material:
  1. Cost is stable, location is not. 49 of 180 route-years fell below 80 percent
     overlap at 1 km, covering 21 of the 32 "robust" links. Pair 17-24 hit zero in 2020.
  2. Vegetation weighting barely changes direction (43 of 45 links keep their sign) but
     strongly changes magnitude (median absolute change 3.6 percent at veg-low, 9.3 at
     veg-high).
  3. The ranking depends on the link graph. At k = 2/3/4/5 you get 31/45/61/77 edges,
     only k = 3 has a bridge, and the top links change at k = 4 and 5.
TRAP: do not present a single ranking as if the graph construction were given.}

\TODO{Outcome of \S\ref{sec:validation}. Report it even if weak.
FACTS: early 2010-2012 = 413 thinned 10 km cells, late 2013-2016 = 450, Jaccard 0.239.
Late cells sit in slightly higher modelled resistance (medians 2.28 vs 2.13, Cliff's
delta 0.11). More late cells fall inside the fixed cores (67.8 vs 45.8 percent).
Almost none fall near a priority path (0 of 413, 3 of 450).
TRAP: this is not independent. The same archive defined the cores, so the core result
is circular by construction. And there are no records after 2016, so 2020 and 2024 have
no check at all. Say both.}

\TODO{H4. No result. The cell-level index and the W1--W3 weight sets were never built.
Fold this into the sensitivity/plausibility subsection: state that the monitoring
output is the link-level robust/uncertain split instead, then explain what was planned,
what was built, and why the gap remains.}

%------------------------------------------------------------------ figure inserts
\begin{figure}[htbp]\centering
  \includegraphics[width=\textwidth]{Figure_2B_current_change_2012_2024.pdf}
  \caption{Modeled current change between 2012 and 2024. The map is a descriptive
  comparison of structural-flow outputs and should not be interpreted as observed
  movement or a validated corridor.}
  \label{fig:change}
\end{figure}

\begin{figure}[htbp]\centering
  \includegraphics[width=\textwidth]{Figure_2A_modelled_current_2024.pdf}
  \caption{Modeled current concentration for 2024, with the fixed core network shown
  for orientation. Current concentration is structural-flow evidence, not occurrence
  probability.}
  \label{fig:current2024}
\end{figure}

\begin{table}[htbp]\centering
  \caption{Covariates, native resolution, source years, and static/dynamic treatment.}
  \label{tab:covariates}
  \small\begin{tabular}{>{\raggedright\arraybackslash}p{0.22\textwidth}>{\raggedright\arraybackslash}p{0.24\textwidth}>{\raggedright\arraybackslash}p{0.13\textwidth}>{\raggedright\arraybackslash}p{0.34\textwidth}}
\toprule
Covariate & Source & Treatment & Temporal role / notes\\
\midrule
Built-up fraction & Global Human Settlement Layer (GHSL) & Dynamic & 2010, 2015, 2020, 2025 products matched to snapshots\\
Nighttime radiance & Visible Infrared Imaging Radiometer Suite (VIIRS), VNL V2 & Dynamic & 2013, 2016, 2020, 2024 products\\
Tree cover & Moderate Resolution Imaging Spectroradiometer (MODIS), MOD44B & Dynamic & Annual snapshot fields\\
Non-tree vegetation & Moderate Resolution Imaging Spectroradiometer (MODIS), MOD44B & Dynamic & Annual snapshot fields\\
Bare cover & Moderate Resolution Imaging Spectroradiometer (MODIS), MOD44B & Dynamic & Annual snapshot fields\\
Road density & Global Roads Inventory Project, version 4 (GRIP4) & Static & Spatial pressure/context\\
Livestock density & Gridded Livestock of the World, version 4 (GLW4) & Static & Combined cattle, goats, and sheep exposure\\
Terrain / slope & Shuttle Radar Topography Mission (SRTM) & Static & Slope-derived resistance\\
Hydrography & Hydrological data and maps based on Shuttle Elevation Derivatives at multiple Scales (HydroSHEDS) & Static & Water/context mask\\
Veterinary fences & World Wide Fund for Nature (WWF), Kavango--Zambezi Transfrontier Conservation Area (KAZA); Kruger references & Scenario & Sensitivity/context, not assumed impassable\\
\bottomrule
\end{tabular}
\\[0.5em]
\parbox{\textwidth}{\footnotesize\emph{Note:} Product shorthand retained in the source names is defined at first use: VNL = VIIRS nighttime lights; V2 = product version 2; MOD44B = MODIS vegetation-continuous-fields product.}

\end{table}

\begin{table}[htbp]\centering
  \caption{Ten highest-ranked modeled priority links shown from the 32-link primary-priority set.}
  \label{tab:top20}
  \small\begin{tabular}{rrrrrrr}
\toprule
From & To & Rank & Betweenness & Type & Distance (km) & Change (\%)\\
\midrule
38 & 40 & 1 & 0.409 & KNN+MST & 120.6 & +4.78\\
24 & 38 & 2 & 0.316 & KNN & 202.3 & +7.44\\
17 & 24 & 3 & 0.300 & KNN+MST & 208.0 & -5.78\\
30 & 40 & 4 & 0.238 & KNN & 169.2 & +3.55\\
27 & 41 & 5 & 0.103 & KNN+MST & 196.8 & -6.45\\
41 & 45 & 6 & 0.090 & KNN & 183.6 & -5.50\\
33 & 38 & 7 & 0.079 & KNN & 108.1 & +13.12\\
9 & 17 & 8 & 0.079 & KNN+MST & 109.0 & -5.26\\
26 & 27 & 9 & 0.079 & KNN+MST & 221.8 & -5.18\\
40 & 43 & 10 & 0.077 & KNN & 159.6 & -1.56\\
\bottomrule
\end{tabular}
\noindent\emph{Note:} The table shows the ten highest-ranked members of the 32-link primary-priority set; links are modeled structural network relationships, not observed animal routes.

\end{table}

\begin{table}[htbp]\centering
  \caption{Protected-area coverage of modeled high-current classes.}
  \label{tab:pa}
  \resizebox{\textwidth}{!}{\begin{tabular}{rrrrrr}
\toprule
Year & Current percentile & Cells & Protected (\%) & Background (\%) & Difference (pp)\\
\midrule
2012 & 80 & 418,295 & 24.23 & 31.26 & -7.03\\
2012 & 90 & 209,148 & 26.31 & 31.26 & -4.95\\
2012 & 95 & 104,574 & 35.57 & 31.26 & +4.31\\
2012 & 99 & 20,915 & 44.17 & 31.26 & +12.92\\
2024 & 80 & 418,292 & 23.27 & 31.26 & -7.99\\
2024 & 90 & 209,146 & 25.54 & 31.26 & -5.72\\
2024 & 95 & 104,573 & 33.45 & 31.26 & +2.19\\
2024 & 99 & 20,915 & 44.63 & 31.26 & +13.37\\
\bottomrule
\end{tabular}}
\\[0.5em]
\noindent\TODO{Final note: values use the outside-core, positive-current domain. Percentage-point differences are relative to the matched background. Do not report a permutation $p$-value because inference was sensitive to the assumed spatial block size.}

\end{table}

\begin{table}[htbp]\centering
  \caption{Overview of completed modeled-connectivity results.}
  \label{tab:priority}
  \begin{tabular}{p{0.53\textwidth}p{0.16\textwidth}p{0.25\textwidth}}
\toprule
Metric & Value & Interpretation\\
\midrule
Final modeled-current snapshots & 4 (2012, 2016, 2020, 2024) & 253/253 focal-region pairs connected each year\\
Fixed density cores & 23 & Historical baseline, fixed across snapshots\\
Selected neighboring links & 45 & Core-pair comparison set\\
Primary-priority links & 32 & Robust under weight/fence screens\\
All-pair median effective-resistance change & -1.224\% & 2012--2024\\
Selected-link median change & -1.557\% & 2012--2024\\
Primary-priority median change & -2.618\% & 2012--2024\\
Largest all-pair increase & +13.119\% & Link 33--38\\
Largest all-pair decrease & -8.243\% & Link 43--45\\
\bottomrule
\end{tabular}

\end{table}
